% translation-source: docs/user/latex/chapters/21-glossaire.tex; sha256: 40ef1adef14fe31a681927ede9ec0e334daefa79c85ba7e25322e6b697e901a1
\chapter{Practical glossary}
\begin{description}
\item[RAW] Shooting data requiring development; rendering depends on decoding,
white balance and profiles, among other factors.
\item[DCP] Camera color profile used in RAW development.
\item[LCP] Lens profile enabling certain optical corrections.
\item[ICC] Colorimetric description of a space or device, distinct from an
interface color palette.
\item[EV] Exposure-change unit; an increase of 1 EV is a linear light gain of
two before output rendering.
\item[Temperature] White-balance setting expressed in kelvins, not a universal
measurement of all image pixels.
\item[Tint] Correction along the green/magenta axis.
\item[Hue] Color family, such as red or orange; distinct from the Tint slider.
\item[Luminance] Relative brightness quantity; precise definition depends on
the space and processing stage.
\item[Gamut] Set of colors representable in a given space.
\item[sRGB] Common output space for the rendered exports described in this book.
\item[Preview] Display image that may have lower resolution than the source;
not proof of export resolution.
\item[Sidecar] File associated with the source, here \chemin{.ipr}, storing
settings without replacing source pixels.
\item[INI] Interface preference file, independent of photo sidecars.
\item[Neutral] Declared adjustment reference for comparison; starting development,
profile and white balance must always be specified.
\end{description}
